\begin{proof}
\begin{enumerate}[label={$\langle1\rangle$\arabic*.},leftmargin=*]
\item 
\noindent\textbf{Step 1. The logarithmic bundle and its restriction.}
The coordinates in \ref{push:src:E.2.17} make $\mathcal W$ a
relative normal-crossing divisor over $s$.
Define
\[
 \mathcal L=
 \ker\left[
 d\log q:
 T_{\mathcal N/\Delta_s}(-\log\mathcal W)
       \longrightarrow\mathcal O_{\mathcal N}
 \right].
 \tag{E.2.19.2}
\]
The map is surjective: if
$q=z_1\cdots z_k$, it sends each
$z_i\partial_{z_i}$, $i\leq k$, to one.
Consequently $\mathcal L$ is locally free of
rank two.  Its sections are actual holomorphic
vector fields, tangent to the branches, and
annihilate both $q$ and $s$.

Let $\nu:\mathcal S\to\mathcal W$ and
$\mathcal D\subset\mathcal S$ be the
simultaneous normalization and conductor.
Restriction to a branch gives a canonical
isomorphism
\[
 \nu^*(\mathcal L|_{\mathcal W})
 \simeq T_{\mathcal S/\Delta_s}(-\log\mathcal D).
 \tag{E.2.19.3}
\]
This can be checked at every stratum.  For
$q=xy$, a basis is
$x\partial_x-y\partial_y,\partial_z$;
on the branch $x=0$ it becomes
$-y\partial_y,\partial_z$, a logarithmic
basis there.  For $q=xyz$, use
$x\partial_x-z\partial_z$ and
$y\partial_y-z\partial_z$.  Their
restrictions to $x=0$ have coefficient
matrix of determinant one in the basis
$y\partial_y,z\partial_z$.
The one-branch case is immediate.
The same computation keeps $s$ fixed.

By \ref{push:src:E.2.18} the right-hand side of
(E.2.19.3) has a global holomorphic
frame.  On the two preimages of each
conductor curve, the identifications
in (E.2.19.3) describe descent by an
invertible two-by-two matrix.  Its
entries are holomorphic on a proper
$\mathbb P^1$ family over $\Delta_s$,
so are constant along each curve
fibre and are holomorphic functions
of $s$.  Denote these matrices by
$G_j(s)$.

For $s\ne0$ both invariant logarithmic
fields in the prescribed Mumford model
descend.  Indeed, its period maps in
multiplicative coordinates multiply
the two coordinates by units depending
only on the base.  Their derivative
on the two vertical logarithmic
fields is the identity.  The toric
small-resolution charts preserve
these same fields.  Thus the entire
two-dimensional space of sections
on the normalized pair descends for
$s\ne0$.  In any chosen frame this
means $G_j(s)=I$, and the identity
theorem gives $G_j(0)=I$ as well.

There is no additional descent
condition hidden at a triple point.
For a normal-crossing local ring,
functions on the normalization
descend precisely when their
restrictions to every pair of
branches agree.  This is the usual
equalizer for
$\C\{x,y,z\}/(xyz)$, and likewise
for one or two branches.  Tensoring
with the locally free bundle
$\mathcal L|_{\mathcal W}$ proves
the same assertion for its sections.
Consequently the chosen normalized
frame descends, and
\[
 \mathcal L|_{\mathcal W}
       \simeq\mathcal O_{\mathcal W}^{\oplus2}.
 \tag{E.2.19.4}
\]

\noindent\textbf{Step 2. A frame on the whole neighbourhood.}
The map $f$ is proper and flat; this
also follows locally from the equations
$(q,s)=(z_1\cdots z_k,s)$.
The bundle $\mathcal L$ is flat over
$B$.  On each ordinary fibre it is
the tangent bundle of a complex
two-torus.  On each $q=0$ fibre it
is trivial by (E.2.19.4).  All these
fibres are compact and connected,
so in every case
\[
 h^0(\mathcal N_{q,s},
       \mathcal L|_{\mathcal N_{q,s}})=2,
\]
and those sections generate the
bundle pointwise.  The proper
constant-cohomology base-change
theorem, in the form established
in \SeriesRef{push:src:E.3.1}, makes $f_*\mathcal L$
a rank-two vector bundle on $B$.
Its evaluation map
\[
 f^*f_*\mathcal L\longrightarrow\mathcal L
 \tag{E.2.19.5}
\]
is an isomorphism on every fibre,
and hence everywhere by Nakayama's
lemma.  Shrink the bidisc and
choose a holomorphic frame $V_1,V_2$.

These fields commute.  On every
ordinary compact torus, every
holomorphic vector field is a
translation field, so their
bracket vanishes there.  The
ordinary locus is dense, and
the bracket is holomorphic.
Each field is complete because
it is vertical and every fibre
of the proper map is compact:
the holomorphic local flow can
be continued for every complex
time on that compact invariant
fibre.  Holomorphic dependence
on initial points and parameters
gives a holomorphic action
\[
 \operatorname{Act}:\C^2\times
       \mathcal N\longrightarrow\mathcal N
\]
over $B$.

\noindent\textbf{Step 3. Holomorphic extension of the two toric periods.}
Choose a smooth point $p_0$ of the
central cusp away from its conductor.
The map $f$ is a submersion there,
so there is a holomorphic section
$p(q,s)$ through $p_0$.
On the central normalized toric
pair the two fields are a basis
of the toric logarithmic Lie
algebra.  Its open orbit is
$(\C^*)^2$.  Let
$v_1^0,v_2^0\in\C^2$ be a
basis of its full period lattice
in the coordinates $V_1,V_2$.
They are complex-linearly
independent.

Near each $v_i^0$, write
$\operatorname{Act}(v,p(q,s))=p(q,s)$
using two vertical coordinates
in a submersion chart around
$p_0$.  The derivative in $v$
is invertible at $(v_i^0,0,0)$,
since the infinitesimal action
is an isomorphism on that open
orbit.  The holomorphic implicit
function theorem gives a unique
holomorphic vector $v_i(q,s)$
near $v_i^0$ solving this equation.
For $q\ne0$, fixing $p(q,s)$
means that this translation is
the identity on the entire
torus fibre.  The equality of
holomorphic maps then extends
over $q=0$.  Thus $v_i(q,s)$
is a period on every fibre.
Their complex-linear
independence persists after
shrinking $B$.

Use these two columns as a
change of the frame.  Because
their coefficients are pulled
back from $B$, the new fields
still commute and are complete.
They have period one, and the
action consequently factors
through
$\C^2/\Z^2\simeq(\C^*)^2$,
using the exponential
$z\mapsto e^{2\pi i z}$
in each coordinate.

\noindent\textbf{Step 4. Primitivity and the actual marking.}
On $B^\circ=\Delta_q^*\times\Delta_s$
the period lattices of the
translation action are the
integral first-homology local
system of the torus family.
The continuously varying
periods $v_1,v_2$ represent
flat integral classes: in a
local differentiable
trivialization their flow
loops vary continuously, and
integral homology is discrete.
They generate a rank-two
invariant sublattice $L$.

Suppose $L$ were not primitive
at some ordinary fibre.  There
would be a rational pair
$(a_1,a_2)\notin\Z^2$ such
that $a_1v_1+a_2v_2$ is an
integral period there.  This
is an equality of flat
integral classes after
multiplication by a common
denominator, so it persists
throughout the connected
ordinary base.  In particular
the corresponding fixed
finite-order element of
$(\C^*)^2$ acts identically
there.  By the identity
theorem it acts identically
on the central open orbit.
This contradicts the choice
of $v_1^0,v_2^0$ as the full
central toric period lattice.
Thus $L$ is primitive.

The marked cusp monodromy,
constant in $s$ by \ref{push:src:E.2.16},
has the form
\[
 \widehat\gamma\longmapsto
       \widehat\gamma-\widehat\delta,
 \qquad
 \widehat u\longmapsto\widehat u+\widehat w,
 \qquad
 \widehat w\longmapsto\widehat w,
 \qquad
 \widehat\delta\longmapsto\widehat\delta.
 \tag{E.2.19.6}
\]
Its full integral invariant
lattice is
$\Z\widehat w\oplus\Z\widehat\delta$.
Since $L$ is a primitive
rank-two sublattice of this
rank-two lattice, equality
follows.  An integral
unimodular change of the two
periods therefore gives
exactly the marking
$\mathbf w,\boldsymbol\delta$.
This also proves that no
extra finite kernel has been
silently introduced.

\noindent\textbf{Step 5. Characters at the triple points.}
The two triple-point sections
are fixed by the torus action:
on the normalization they
are toric vertices.  Restrict
the derivative representation
to the compact real torus
$(S^1)^2$ and average a
Hermitian metric along either
section.  Its characters
are integral pairs and
vary continuously in $s$;
hence their multiset is
locally constant.  Equivalently,
Fourier averaging gives the
individual character
subbundles.  They coincide
at $s=0$ with the known
characters for $s\ne0$.

For clarity, the two unit
triangles in the prescribed
cocharacter lattice have
height-one rays
\[
 \{(0,0,1),(1,0,1),(0,1,1)\},
 \qquad
 \{(1,0,1),(0,1,1),(1,1,1)\}.
\]
Inverting their unimodular
ray matrices gives the two
character triples in
(E.2.19.1); the sum in
each triple is zero, as
required by invariance of
$q=z_1z_2z_3$.
The calculation is recorded
in \path{code/d_push_cusp_torus_action.py}.
All action and marking
assertions follow.

\end{enumerate}
\end{proof}
